\chapter{Tests}\label{chap:tests}

\section{Stratégie}
Trois niveaux de vérification se complètent. Les contrôles automatisés couvrent le métier, les utilitaires, la
persistance, les parties pures de la couche web et la structure statique des pages ; ils s'exécutent par une seule
commande et affichent un total. Le parcours manuel vérifie ce qu'aucun test sans serveur ne peut voir (rendu,
cycle de vie JSF, redirections). Les captures de ce rapport en sont la trace.

\section{Contrôles automatisés}
La classe \code{Verifications} enchaîne six lots, sans bibliothèque de test : une petite classe \code{Checks} compte les
réussites et les échecs. Le résultat attendu est \nbChecks{} réussites et aucun échec.

\begin{table}[htbp]
\centering
\caption{Lots de contrôles}\label{tab:lots}
\begin{tabularx}{\linewidth}{@{}lX@{}}
\toprule
\textbf{Lot} & \textbf{Ce qu'il vérifie}\\
\midrule
\code{PanierCheck} & Règles R1 à R6 du panier : bornes de quantité, stock, fusion de lignes, total en \code{BigDecimal}.\\
\code{ServiceChecks} & Services métier : validation des formulaires, messages d'erreur.\\
\code{UtilChecks} & Hachage PBKDF2, argent, limiteur de tentatives.\\
\code{JpaChecks} & Persistance réelle sur MySQL : contraintes refusées, verrouillage, héritage (nécessite la base).\\
\code{WebChecks} & Parties pures de la couche web : redirections sûres, jeton CSRF, politique CSP, adresses de la vitrine, aides d'affichage, page d'erreur.\\
\code{VuesChecks} & Analyse statique des pages Facelets (section~\ref{sec:vues-checks}).\\
\bottomrule
\end{tabularx}
\end{table}

\section{Tests de la couche web}\label{sec:webchecks}
\code{WebChecks} montre l'intérêt de sortir la logique des pages et des beans dans des méthodes pures. Quelques exemples :

\begin{itemize}
  \item \code{Chemins.cheminSur} refuse \code{//site-externe}, les URL absolues, l'antislash et un saut de ligne
        (injection d'en-tête) ; le chemin interne est accepté.
  \item \code{VitrineBean.construireUrl} encode la recherche en UTF-8 (\code{café noir} $\rightarrow$ \code{caf\%C3\%A9+noir}),
        omet la page~1 et neutralise un \code{\&page=9} saisi dans la recherche.
  \item \code{SecurityFilter.egal} compare les jetons sans fuite par le temps et refuse un jeton absent ou de longueur
        différente.
  \item \code{Affichage.prix} formate avec deux décimales, une virgule et une espace insécable avant « MAD ».
\end{itemize}

\section{Analyse statique des vues}\label{sec:vues-checks}
Une erreur dans une page Facelets n'apparaît normalement qu'au chargement de la page. \code{VuesChecks} la détecte avant, sans serveur :

\begin{enumerate}
  \item chaque page, gabarit et composant est un document XML bien formé ;
  \item aucun \code{h:commandLink}, aucun \code{f:ajax}, aucun attribut \code{on*} (cohérence avec la CSP) ;
  \item chaque \code{h:form} contient \code{t:csrf} ;
  \item chaque expression \code{bean.membre} désigne un bean existant et un accesseur ou une méthode qui existe.
\end{enumerate}

La dernière règle remplace une partie de ce qu'un EL mal écrit ferait découvrir en pleine démonstration. Elle a déjà trouvé
un faux positif instructif : une expression EL laissée dans un commentaire XML d'un composant. Facelets évalue l'EL même dans les
commentaires ; le commentaire a été reformulé.

\section{Parcours manuel}
Le tableau~\ref{tab:parcours} liste les quinze scénarios déroulés dans le navigateur. Chaque capture porte la barre d'adresse
et le titre de l'onglet, afin que la page ouverte soit identifiable.

\begin{table}[htbp]
\centering
\caption{Scénarios du parcours manuel}\label{tab:parcours}
\small
\begin{tabularx}{\linewidth}{@{}rXl@{}}
\toprule
\textbf{N°} & \textbf{Scénario et résultat attendu} & \textbf{Figure}\\
\midrule
1 & Ouvrir la vitrine : cartes, pastilles de catégories. & \ref{fig:vitrine}\\
2 & Rechercher, puis filtrer : l'adresse porte \code{q} et \code{categorie}, le compteur change. & \ref{fig:recherche}\\
3 & Ouvrir un produit : prix, stock, « Se connecter pour acheter ». & \ref{fig:detail}\\
4 & Connexion avec un mauvais mot de passe : message rouge, champ vidé. & \ref{fig:connexion}\\
5 & Connexion correcte : « Bienvenue », badge du panier. & ---\\
6 & Ouvrir le panier : lignes, quantités, total. & \ref{fig:panier}\\
7 & Dépasser le stock du produit \code{INF-003} : refus, panier inchangé ; produit à stock nul : bouton désactivé. & \ref{fig:stock}\\
8 & Ajouter un produit déjà présent : fusion de la ligne ; \code{F5} ne rejoue pas le formulaire. & ---\\
9 & Contrôler la base : les lignes correspondent à l'écran. & \ref{fig:controle}\\
10 & Se déconnecter : retour à la vitrine avec message. & ---\\
11 & Connexion administrateur : arrivée sur la liste des produits. & \ref{fig:admin}\\
12 & Nouveau produit, formulaire vide : erreurs sous les champs. & \ref{fig:adminform}\\
13 & Inscription avec un mot de passe trop court : erreurs, mots de passe vidés. & \ref{fig:inscription}\\
14 & Client devant \code{admin/produits.xhtml} : page 403. & \ref{fig:403}\\
15 & Non connecté devant \code{panier.xhtml} : redirection vers la connexion, puis retour au panier. & ---\\
\bottomrule
\end{tabularx}
\end{table}

Le scénario~8 vérifie le motif Post/Redirect/Get : après l'ajout, la page affichée est le résultat d'une redirection
(\code{GET}), donc un rechargement n'envoie rien. Le scénario~7 se fait par ajouts répétés : un navigateur bloque
de toute façon une saisie au-dessus du maximum d'un champ numérique, c'est pourquoi la règle est aussi vérifiée côté serveur.

\fig{13-mysql-controle-paniers}{Contrôle en base : lignes du panier après les manipulations}{fig:controle}
